% Supplementary worked example. Copy the subsection into the supplement.
% Figure numbering follows the user's Extended Data Figure 4.
\documentclass[11pt,a4paper]{article}
\usepackage[margin=22mm]{geometry}
\usepackage[T1]{fontenc}
\usepackage{lmodern}
\usepackage{amsmath,booktabs}
\setlength{\parindent}{0pt}
\setlength{\parskip}{5pt}
\begin{document}

\subsection*{Calculating \(N_{\mathrm{eff}}^{+}\) and \(P_c\): a worked example}

Consider a time slice crossing the three terminal branches in each scenario of Extended Data Figure~4. Code each lineage in the order \emph{wings, thorax, head}, with 1 for accelerated and 0 otherwise. Each complete combination is one pattern: \texttt{111} is a single pattern even though all three units accelerate.

\paragraph{1. Calculating phenomic diversity, \(N_{\mathrm{eff}}^{+}\).}
At time slice \(t\), let \(N_t\) be the number of branches crossing the slice and \(N_t^{+}\) the number containing at least one accelerated unit. Among these \(N_t^{+}\) branches, calculate the relative frequency \(p_i(t)\) of each distinct acceleration pattern. Using natural logarithms,
\[
f^{+}(t)=\frac{N_t^{+}}{N_t},\qquad
H^{+}(t)=-\sum_i p_i(t)\ln p_i(t),\qquad
N_{\mathrm{eff}}^{+}(t)=f^{+}(t)e^{H^{+}(t)}.
\]
All-zero patterns are excluded from \(H^{+}\), but their branches remain in the denominator of \(f^{+}\).

In Scenario~3, the patterns \texttt{001}, \texttt{010} and \texttt{100} each occur once. Thus \(f^{+}(t)=3/3=1\) and each \(p_i(t)=1/3\), giving
\[
H^{+}(t)=-3\left(\frac13\ln\frac13\right)=\ln3,
\qquad N_{\mathrm{eff}}^{+}(t)=1\times e^{\ln3}=3.
\]
If no branch contains an accelerated unit (Scenario~4), define \(H^{+}(t)=0\) and \(N_{\mathrm{eff}}^{+}(t)=0\).

\paragraph{2. Calculate phenomic coordination, \(P_c\).}
For each branch, divide its number of accelerated units, \(a_b\), by the total number of units, \(L\), then average across \emph{all} branches crossing time slice \(t\). In Scenario~3, each branch has one accelerated unit out of three:
\[
P_c(t)=\frac{1}{N_t}\sum_{b=1}^{N_t}\frac{a_b}{L}
=\frac{1/3+1/3+1/3}{3}=\frac13.
\]

\begin{center}
\begin{tabular}{clcccc}
\toprule
Scenario & Lineage patterns & \(f^{+}\) & \(H^{+}\) & \(N_{\mathrm{eff}}^{+}\) & \(P_c\) \\
\midrule
1 & \texttt{111, 111, 111} & 1 & 0 & 1 & 1 \\
2 & \texttt{100, 100, 100} & 1 & 0 & 1 & \(1/3\) \\
3 & \texttt{001, 010, 100} & 1 & \(\ln3\) & 3 & \(1/3\) \\
4 & \texttt{000, 000, 000} & 0 & 0 & 0 & 0 \\
\bottomrule
\end{tabular}
\end{center}

Scenario~1 has the highest \(P_c\) because all units accelerate on every lineage. Scenario~3 has the highest \(N_{\mathrm{eff}}^{+}\) because different units accelerate on different lineages. The two metrics therefore describe different properties; \(P_c\) cannot be obtained from the \(N_{\mathrm{eff}}^{+}\) value alone.

To illustrate the effect of an unaccelerated lineage, replace \texttt{100} in Scenario~3 with \texttt{000}. Then \(f^{+}=2/3\), the two remaining patterns each have frequency \(1/2\), and \(H^{+}=\ln2\), giving \(N_{\mathrm{eff}}^{+}=4/3\) and \(P_c=2/9\). The zero pattern is excluded from \(H^{+}\), but its lineage remains in the denominators of \(f^{+}\) and \(P_c\).

\end{document}
